Una partícula se mueve según $\vec{r}(t) = (2\sin t,\; 2\cos t,\; 3t)$ (distancias en metros).
\begin{enumerate}
    \item Calcular la velocidad $\vec{v}(t)$, la aceleración $\vec{a}(t)$ y la rapidez $v(t)$.
    \item Descomponer la aceleración en sus componentes tangencial y normal:
    \[ a_T = \frac{\vec{r}'(t)\cdot\vec{r}''(t)}{\|\vec{r}'(t)\|}, \qquad a_N = \frac{\|\vec{r}'(t)\times\vec{r}''(t)\|}{\|\vec{r}'(t)\|} \]
    \item Verificar que $\|\vec{a}\|^2 = a_T^2 + a_N^2$.
    \item Calcular la curvatura $\kappa$ de la trayectoria. ¿Cuál es su radio de curvatura $\rho = 1/\kappa$?
\end{enumerate}